\chapter{Classification}
\label{ch:classification}

The operational states of Chapter~\ref{ch:states} were defined from motion alone. This
chapter asks the leakage-free question that motivates the whole design: can the
\emph{air-system} sensors, which never touched that definition, recover the deceleration
state on their own? If they can, deceleration leaves an independent fingerprint across the
pneumatic system, which is the premise the later analysis rests on.

\section{Setup}

The target is the binary label \enquote{deceleration versus not} on ten-second windows.
The motion channels that defined the target are excluded from the predictors entirely, in
line with the role ledger (Section~\ref{sec:roles}). Two predictor tiers are compared:

\begin{itemize}
  \item \textbf{auxiliary only} (60 features): the channels classed \emph{auxiliary} in the
        role ledger (Section~\ref{sec:roles}): the health and support channels (main-reservoir
        and load pressure, load signal, energy-braking resistance, ambient temperature, and
        compressor state). They are termed auxiliary because they never carry the brake command:
        they report the state of the air system rather than the braking demand itself, so
        recovering deceleration from them alone is the strictest test of independent signal.
  \item \textbf{full} (180 features): the auxiliary channels together with the \emph{actuation}
        channels (the brake command and its immediate pneumatic response), which, unlike the
        auxiliary channels, are driven directly by the braking demand.
\end{itemize}

Several transparent models are fitted so that the result does not hinge on one method's
quirks: a decision tree, LDA, QDA, logistic regression, and a random forest. All are
trained on the training window and scored on the held-out window
(Section~\ref{sec:split}). Because deceleration is the minority class (about 20\% of
windows), balanced accuracy and ROC-AUC are read in preference to raw accuracy.

\section{Results}

The headline numbers are in Table~\ref{tab:clf}. The full-tier random forest reaches a
ROC-AUC of 0.954 on held-out data, with balanced accuracy 0.869 and recall 0.956 for the
deceleration class. Its precision is lower, 0.528: the model catches almost all
decelerations but also flags a fair number of non-decelerations, which is the expected
trade-off when recall is high on a minority class and is the reason balanced accuracy and
AUC are the yardsticks here rather than precision alone. LDA and the decision tree follow
closely (AUC 0.924 and 0.910); the decision tree is the simplest model yet lands near the
top, making it the most interpretable strong option. QDA reaches 0.901 and logistic
regression 0.908. Full per-model metrics for both tiers are collected in
Appendix~\ref{app:clfmetrics}.

The auxiliary-only tier result is the most interesting. With the brake command \emph{withheld}, the
random forest still reaches AUC 0.921. Deceleration therefore leaves a fingerprint across
the independent health sensors (reservoir pressure, load pressure, and the rest) even
without being told the brake was applied. Most of the recoverable signal lives in the
auxiliary channels, not only in the actuation channels, which is what makes the pneumatic
system a plausible place to look for wear.

\subsection*{Which channels carry the signal}
\label{sec:clf-importance}

The interpretable models agree on where the deceleration signal sits. The decision tree
puts almost all of its weight on a single channel, pneumatic braking force, with spring-brake
pressure and proportional-valve pressure filling in the lower splits
(Table~\ref{tab:clfimp}). LDA spreads its weight more evenly but ranks the same subsystem
first: the per-wagon pneumatic-braking-force and proportional-valve-pressure copies dominate
its discriminant, with reservoir pressure the strongest auxiliary contributor. QDA exposes
no single loading vector, but it reaches comparable AUC on the same channel pool, so all
three lean on the same part of the pneumatic system. The channels that matter are the ones
that actuate and support braking, and none of them entered the motion-only state definition.

\begin{table}[htbp]
  \centering
  \caption{Top air-system channels for the full-tier classifiers. The decision-tree column
  is the split importance; the LDA column is the standardised discriminant weight
  ($|\text{coef}|$). Per-wagon copies are collapsed to their channel family.}
  \label{tab:clfimp}
  \begin{tabular}{lrr}
    \toprule
    Channel (family) & Decision tree & LDA $|\text{coef}|$ \\
    \midrule
    pneumatic braking force        & 0.69 & high \\
    spring-brake pressure          & 0.14 & low \\
    proportional-valve pressure    & 0.11 & highest \\
    energy-braking resistance      & 0.05 & low \\
    main-reservoir pressure        & 0.00 & moderate \\
    \bottomrule
  \end{tabular}
\end{table}

\begin{table}[htbp]
  \centering
  \caption{Deceleration classification on the held-out window. The full-tier random forest
  leads; the auxiliary-only random forest, with the brake command withheld, is nearly as
  strong. AUC is measured against 0.5 (coin flip); balanced accuracy against per-class
  chance.}
  \label{tab:clf}
  \begin{tabular}{llrrr}
    \toprule
    Tier & Model & ROC-AUC & Balanced acc. & $F_1$ (decel.) \\
    \midrule
    full          & random forest & 0.954 & 0.869 & 0.680 \\
    full          & LDA           & 0.924 & 0.726 & 0.610 \\
    full          & decision tree & 0.910 & 0.825 & 0.609 \\
    full          & QDA           & 0.901 & 0.746 & 0.602 \\
    auxiliary only & random forest & 0.921 & 0.831 & 0.620 \\
    \bottomrule
  \end{tabular}
\end{table}

\begin{figure}[htbp]
  \centering
  \includegraphics[width=0.9\linewidth]{state_clf_comparison.png}
  \caption{Model comparison for deceleration recovery (held-out AUC and $F_1$). The
  full-tier random forest leads; the auxiliary-only tier, without the brake command,
  remains close.}
  \label{fig:clfcomparison}
\end{figure}

\begin{figure}[htbp]
  \centering
  \includegraphics[width=0.85\linewidth]{state_clf_confusion.png}
  \caption{The four most informative confusion matrices on the held-out window
  (row-normalised). The full-tier random forest catches almost every deceleration at the
  cost of some false positives; LDA is the opposite, precise but conservative; the decision
  tree sits between them; and the auxiliary-only random forest, with the brake command
  withheld, keeps the same shape, showing the signal survives without actuation channels.}
  \label{fig:confusion}
\end{figure}
